% !TeX root = ../traffic_rules.tex

\section{Traffic Rule Formalization}
\label{sec:rule_formalization}

\subsection{Traffic Law Terminology}
\label{sec:def_law_terms}
\begin{itemize}
	\item access ramp / acceleration lane
	\item off ramp / deceleration lane
	\item main carriage way
	\item motor road / motor street / highway
	\item ego vehicle
	\item other vehicle
\end{itemize}

\subsection{General Traffic Rules}
\label{sec:general_rules}

\subsubsection{Distance}
\begin{table}[H]
\centering
\begin{tabular}[t]{ll}
\toprule
\textbf{Safe distance front} & \\
\midrule
StVO \S 4 (1)& Der Abstand zu einem vorausfahrenden Fahrzeug muss in der Regel \\ 
& so groß sein, dass auch dann hinter diesem gehalten werden kann, wenn es plötzlich  \\
& gebremst wird.\\

StVO English & A person operating a vehicle moving behind another vehicle must, as a rule,\\
& keep a sufficient distance from that other vehicle so as to be able to pull up safely \\
& even if it suddenly slows down or stops.\\

VCoRT \S 13 (5) & The  driver  of  a  vehicle  moving  behind  another  vehicle  shall  keep  at \\ 
& a  sufficient distance from that other vehicle to avoid collision if the vehicle in front \\ 
& should suddenly slow down or stop. \\

Concretization &  A vehicle moving behind other vehicles within a lane must keep a \\ 
& distance which ensures collision freedom to each of these other vehicles even if one\\ 
&  or several of them are suddenly slow down or stop.\\

%Atomic Prepositions & $\mathtt{keeps\_safe\_distance}$\\

Predicates & $\safeDistanceFrontPredicate$\\ 

Comments &  Evaluated for all vehicles for which the following condition is valid: \\
& $\onSameLaneAsEgoVehiclePredicate \land \frontOfEgoVehiclePredicate$\\

LTL & G$(\safeDistanceFrontPredicate)$\\
MTL & G$(\safeDistanceFrontPredicate)$\\
STL & G$(\safeDistance(\vEgo, \vOther) \leq (\rearPositionFunction(\otherVehicle) - \frontPositionFunction(\egoVehicle)))$\\
\bottomrule
\end{tabular}
\end{table}

\subsubsection{Braking}
\begin{table}[H]
\centering
\begin{tabular}[t]{ll}
\toprule
\textbf{Unnecessary braking} & \\
\midrule
StVO \S 4 (1) & Wer vorausfährt, darf nicht ohne zwingenden
Grund stark bremsen.\\
 & \\
StVO English & The person operating the vehicle in front must not brake suddenly \\
& without a compelling reason.\\


VCoRT \S 17 (1) &  No driver of a vehicle shall brake abruptly unless it is necessary to do so for\\
&  safety reasons.\\

Concretization &  A vehicle shall never brake abruptly unless it is necessary to do so \\
& for safety reasons. A vehicle brakes abruptly if: 1. It brakes with an acceleration \\ 
& which violates an user-defined threshold and there is no other vehicle, obstacle, \\ 
& or speed limit front of the vehicle; 2. There is an obstacle in front of the vehicle \\ 
& and the acceleration difference and jerk difference between the vehicle and the obstacle \\ 
& violates an acceleration and jerk threshold, respectively; 3. There is an speed limit \\
& in front of the vehicle and the acceleration and jerk of the vehicle violate \\
& an acceleration and jerk threshold, respectively.\\

Predicates & $\unnecessaryBraking$\\

Comments & \\

LTL & G$(\unnecessaryBraking)$\\
MTL & G$(\unnecessaryBraking)$\\
STL & \\
\bottomrule
\end{tabular}
\end{table}%

\subsubsection{Velocity}
\begin{table}[H]
\centering
\begin{tabular}[t]{ll}
\toprule
\textbf{Maximum Speed Limit} &\\
\midrule
StVO \S 3 (1), \S 3 (3),  & The legal text is omitted for reasons of space.\\
\S 18 (1), \S 18 (5), \S 18 (6) & \\
StVO English & The legal text is omitted for reasons of space.\\

VCoRT \S 13 (2) &  Domestic legislation shall establish maximum speed limits for all roads. \\
& Domestic legislation shall also determine special speed limits applicable \\ 
& to certain categories of vehicles presenting a special danger, in particular by reason \\ & of their mass or their  load. They may establish similar provisions \\
& for certain categories of drivers, in particular for new drivers.\\

Concretization &  A vehicle shall not exceed the allowed speed limit of the lanes it drives on, shall \\
& not exceed the speed limit which is necessary to ensure collision freedom with \\
& respect to vehicles outside the field of view, and shall not exceed as speed which \\
& is necessary to comfortably ensure traffic regulations.\\

Predicates & $\laneSpeedLimitPredicate$, $\fovSpeedLimitPredicate$, $\brakingSpeedLimitPredicate$\\

Comments &  The traffic laws do not provide concrete information about speed limit, \\ 
& e.g., the VCoRT passes the task of defining speed limits to national lawmakers. \\
& The concretization assumes that a mapping of speed limits provided \\ 
& by traffic signs to lanes is given. If no speed limit is given, \\ 
& the speed limit caused by the field of view (traffic sign based speed limit set \\ 
& to infinity) and a speed limit which ensures comfortable braking for traffic regulations \\
& (e.g., upcoming speed limit) must be considered. The velocity with respect to vehicles \\
& inside the field of view is considered in the safe distance traffic rule concretization.\\

LTL & G$(\laneSpeedLimitPredicate \land \fovSpeedLimitPredicate \land \brakingSpeedLimitPredicate)$\\
MTL & G$(\laneSpeedLimitPredicate \land \fovSpeedLimitPredicate \land \brakingSpeedLimitPredicate)$\\
STL & G($\vEgo \leq \speedLimitLane \land \vEgo \leq \speedLimitFOV)$\\ 
\bottomrule
\end{tabular}
\end{table}

\begin{table}[H]
	\centering
	\begin{tabular}[t]{ll}
		\toprule
		\textbf{Minimum Speed Limit} &\\
		\midrule
		StVO \S 3 (2), sign $275$ & Ohne triftigen Grund dürfen Kraftfahrzeuge nicht so langsam fahren, dass sie\\
		& den Verkehrsfluss behindern.\\
		& \\
		& Zeichen 275: Wer ein Fahrzeug führt, darf nicht langsamer als mit der angegebenen \\ 
		& Mindestgeschwindigkeit fahren, sofern nicht Straßen-, Verkehrs-, Sicht- oder \\ 
		& Wetterverhältnisse dazu verpflichten. Es verbietet, mit Fahrzeugen, die nicht so \\
		& schnell fahren können oder dürfen, einen so gekennzeichneten Fahrstreifen zu benutzen. \\
		
		StVO English & No motor vehicle must, without good reason, travel so slowly as to impede\\ 
		& the flow of traffic.\\
		& \\
		& A person operating a vehicle must not do so below the indicated minimum speed \\ 
		& unless road, traffic, visibility or weather conditions make this necessary.  \\
		& This sign prohibits vehicles that are unable or not permitted to travel as fast from \\
		& using a lane marked in this manner. \\
		
		VCoRT \S 13 (2) &  Similar to maximum speed limit.\\
		
		Concretization &  The relative velocity of a vehicle with respect to leading and following vehicles should \\
		& not exceed a certain threshold.\\
		
		Predicates & $\minSpeedLimitPredicate$\\
		
		Comments &  \\
		
		LTL & G$(\minSpeedLimitPredicate)$\\
		MTL & G$(\minSpeedLimitPredicate)$\\
		STL & G($\vOther^\mathrm{min} - \vEgo \leq v_\mathrm{min}^\mathrm{slow})$\\
		\bottomrule
	\end{tabular}
\end{table}




\subsection{Highway Traffic Rules}
\label{sec:highway_rules}
\begin{table}[H]
	\centering
	\begin{tabular}[t]{ll}
		\toprule
		\textbf{Driving right} & \\
		\textbf{faster than left} &\\
		\midrule
		StVO \S 7 (2), \S 7 (2a),  & Ist der Verkehr so dicht, dass sich auf den Fahrstreifen für eine Richtung \\
		\S 7a & Fahrzeugschlangen gebildet haben, darf rechts schneller als links gefahren werden.\\
		& \\
		& Wenn auf der Fahrbahn für eine Richtung eine Fahrzeugschlange auf dem jeweils linken \\ 
		& Fahrstreifen steht oder langsam fährt, dürfen Fahrzeuge diese mit geringfügig höherer \\ 
		& Geschwindigkeit und mit äußerster Vorsicht rechts überholen. \\ 
		& \\
		& Auf Ausfädelungsstreifen darf nicht schneller gefahren werden als auf den \\  
		& durchgehenden Fahrstreifen. Stockt oder steht der Verkehr auf den durchgehenden \\ 
		& Fahrstreifen, darf auf dem Ausfädelungsstreifen mit mäßiger Geschwindigkeit  \\
		& und besonderer Vorsicht überholt werden. \\
		& \\
		& Gehen Fahrstreifen, insbesondere auf Autobahnen und Kraftfahrstraßen, \\
		& von der durchgehenden Fahrbahn ab, darf beim Abbiegen vom Beginn einer breiten \\
		& Leitlinie (Zeichen 340) rechts von dieser schneller als auf der durchgehenden Fahrbahn \\
		& gefahren werden. \\ 
		& \\
		& Auf Ausfädelungsstreifen darf nicht schneller gefahren werden als auf den durchgehenden \\ 
		& Fahrstreifen. Stockt oder steht der Verkehr auf den durchgehenden Fahrstreifen, darf \\ 
		& auf dem Ausfädelungsstreifen mit mäßiger Geschwindigkeit und besonderer Vorsicht \\
		& überholt werden. \\
		
		StVO English & If the density of traffic has resulted in queues of vehicles in the lanes of one \\
		& carriageway, traffic on the right (nearside lane, middle lane) may move faster than \\
		& traffic on the left (offside lane, middle lane).\\
		& \\
		& If, on a carriageway for one direction of traffic, a vehicle queue has formed and come to \\ 
		& a standstill or is moving at low speed in the left-hand lane, vehicles may overtake this \\ 
		& queue on the right at a slightly higher speed and with the utmost care. \\
		& \\
		& In deceleration lanes, vehicles must not travel faster than traffic on the main carriageway. \\
		& If traffic on the main carriageway is moving slowly or is stationary, vehicles in a \\
		& deceleration lane may overtake at a moderate speed and with great care. \\
		& \\
		& Where lanes diverge from the main carriageway, especially on motorways and motor roads, \\
		& vehicles turning off may, from the beginning of and to the right of a lane marking with \\
		& broad lines (sign 340), travel faster than traffic on the main carriageway. \\
		& \\
		& On motorways and other roads outside built-up areas, vehicles may travel faster in \\
		& acceleration lanes than traffic on the main carriageway. \\
	\end{tabular}
\end{table}

\begin{table}[H]
	\centering
	\begin{tabular}[t]{ll}
		\toprule
		\textbf{Driving right} & \\
		\textbf{faster than left} &\\
		\midrule
		VCoRT  &  \\
		
		Concretization &  A vehicle shall not drive faster than any vehicle on its left lanes. \\
		& A exception is, if $1.$ the vehicle on the left lane is in a congestion; \\
		& $2.$ the vehicle on the left lane is on the main carriage way and the vehicle on the right \\
		& lane is on a off ramp; $3.$ the left vehicle and right vehicle are separated by a broad \\
		& lane marking;  $4.$ the right vehicle is within an access ramp. For $1.$ and $2.$, the vehicle \\ 
		& on the right lane is allowed to drive faster than the vehicle on the left lane \\
		& with a slightly higher speed.\\
		
		Predicates & $\drivesFasterThanVehicleLeftPredicate$, $\congestionLeftPredicate$, $\otherVehicleLeftOfBroadLaneMarkingPredicate$,\\
		& $\drivesMaxWithSlightlyHigherSpeedPredicate$, $\rightOfBroadLaneMarkingPredicate$\\
		
		Comments &  Assumption: Deceleration lane always right of main carriage way.\\
		& Case two can be neglected -> same as case one \\
		
		LTL & G$(\neg \drivesFasterThanVehicleLeftPredicate$ \\ 
		& $\lor (\congestionLeftPredicate \land \drivesMaxWithSlightlyHigherSpeedPredicate)$\\
		& $\lor (\rightOfBroadLaneMarkingPredicate \land \otherVehicleLeftOfBroadLaneMarkingPredicate)$\\
		& $\lor (\onAccessRampPredicate \land \neg \congestionLeftPredicate))$\\
		MTL & G$(\neg \drivesFasterThanVehicleLeftPredicate$ \\ 
		& $\lor (\congestionLeftPredicate \land \drivesMaxWithSlightlyHigherSpeedPredicate))$\\
		STL & \\
		\bottomrule
	\end{tabular}
\end{table}


\begin{table}[H]
	\centering
	\begin{tabular}[t]{ll}
		\toprule
		\textbf{Emergency Lane} &\\
		\midrule
		StVO \S 3 (2) & Sobald Fahrzeuge auf Autobahnen sowie auf Außerortsstraßen mit mindestens zwei \\
		& Fahrstreifen für eine Richtung mit Schrittgeschwindigkeit fahren oder sich die Fahrzeuge \\
		& im Stillstand befinden, müssen diese Fahrzeuge für die Durchfahrt von Polizei- \\
		& und Hilfsfahrzeugen zwischen dem äußerst linken und dem unmittelbar rechts \\
		& daneben liegenden Fahrstreifen für eine Richtung eine freie Gasse bilden.\\
		
		StVO English & As soon as vehicles on motorways and roads outside a built-up area with at least two lanes \\
		& for one direction start to move at walking pace or come to a standstill, these vehicles must \\
		& leave a gap for one direction between the lane on the far left and the lane immediately \\
		& adjacent to it on the right to allow police and emergency vehicles to pass.\\
		
		VCoRT  &  \\
		
		Concretization &  A vehicle in a congestion outside urban roads has to drive leftmost if it is in the leftmost \\
		& lane and has to drive rightmost otherwise. \\
		& not exceed a certain threshold.\\
		
		Predicates & $\minSpeedLimitPredicate$\\
		
		Comments &  \\
		
		LTL & G$(\minSpeedLimitPredicate)$\\
		MTL & G$(\minSpeedLimitPredicate)$\\
		STL & G($\vOther^\mathrm{min} - \vEgo \leq v_\mathrm{min}^\mathrm{slow})$\\
		\bottomrule
	\end{tabular}
\end{table}


\subsection{Urban Traffic Rules}
\label{sec:urban_rules}

\subsection{Cross-Country Road Traffic Rules}
\label{sec:cross_country_rules}